% Part: sets-functions-relations
% Chapter: functions
% Section: composition

\documentclass[../../../include/open-logic-section]{subfiles}

\begin{document}

\olfileid{sfr}{fun}{cmp}
\olsection{Komposisi Fungsi}

\begin{explain}
\oliflabeldef{sfr:fun:inv:sec}{Ing \olref[inv]{sec} wis kita deleng yen
invers~$f^{-1}$ saka !!a{bijection}~$f$ iku uga fungsi. Operasi
liyane ing fungsi yaiku komposisi: k}{K}ita bisa netepake fungsi
anyar kanthi nyusun rong fungsi, $f$ lan~$g$, yaiku kanthi ngetrapake
$f$ dhisik banjur~$g$. Mesthi wae, iki mung bisa ditindakake yen
jangkauan lan domaine trep, yaiku jangkauan saka~$f$ kudu subhimpunan
saka domain~$g$. \oliflabeldef{sfr:rel:ops:sec}{Operasi ing
fungsi iki salaras karo operasi produk relatif ing
relasi saka \olref[rel][ops]{sec}.}{}

Diagram bisa mbantu njlentrehake gagasan komposisi. Ing
\olref{fig:composition}, kita nggambarake rong fungsi $f \colon A \to B$
lan $g \colon B \to C$ sarta komposisine~$(\comp{f}{g})$. Fungsi
$(\comp{f}{g}) \colon A \to C$ masangake saben !!{element} saka~$A$
karo !!a{element} saka~$C$. Kanggo nemtokake !!{element} saka~$C$
kang dadi pasangane !!a{element} saka $A$, carane mangkene: yen diwenehi masukan $x \in
A$, fungsi $f$ ditrapake dhisik tumrap~$x$, kang ngasilake $f(x)
= y \in B$, banjur fungsi $g$ ditrapake tumrap~$y$, kang ngasilake
$g(f(x)) = g(y) = z \in C$.
\begin{figure}
  \olasset[2\olphotowidth]{assets/diagrams/composition.tikz}
  \caption{Komposisi $g \circ f$ saka rong fungsi $f$ lan~$g$.}
  \ollabel{fig:composition}
\end{figure}
\end{explain}

\begin{defn}[Komposisi]
Upamakna $f\colon A \to B$ lan $g\colon B \to C$ iku fungsi.
\emph{Komposisi} saka $f$ karo~$g$ yaiku $\comp{f}{g} \colon A \to C$,
kanthi $(\comp{f}{g})(x) = g(f(x))$.
\end{defn}

\begin{ex}
Gatekna fungsi $f(x) = x + 1$ lan $g(x) = 2x$. Amarga
$(\comp{f}{g})(x) = g(f(x))$, kanggo saben masukan~$x$ kita kudu njupuk
paneruse dhisik, banjur asil mau dipingake loro. Dadi, komposisine
ditetepake kanthi $(\comp{f}{g})(x) = 2(x+1)$.
\end{ex}

\begin{prob}
Tuduhna yen $f \colon A \to B$ lan $g \colon B \to C$ loro-lorone
!!{injective}, mula $\comp{f}{g}\colon A \to C$ iku !!{injective}.
\end{prob}

\begin{prob}
Tuduhna yen $f \colon A \to B$ lan $g \colon B \to C$ loro-lorone
!!{surjective}, mula $\comp{f}{g}\colon A \to C$ iku !!{surjective}.
\end{prob}

\begin{prob}
Upamakna $f \colon A \to B$ lan $g \colon B \to C$. Tuduhna yen graf
saka $\comp{f}{g}$ yaiku $R_f \mid R_g$.
\end{prob}

\end{document}
